\scanpage{49}

d'autres systèmes de gén.

\[
\Sigma_{r'}=(\rho_f^{-1},\rho_s^{-1},\rho_a=\rho_{\unsure{a|r}}^{-1})
\]

associés aux repères $r'=(s,a,f)$ \uline{indirects}
Les $\Sigma_{r'}$ sont conjugués entre eux sous $H$, et les $\Sigma_r$ et $\Sigma_{r'}$ sont conjugués sous $\mathrm{Aut}\,\Pi$, mais pas néc. sous $H=\mathrm{Aut}^{+}(\Pi)$ -- par exemple parce qu'il se peut que les $\rho_s$ et $\rho_s^{-1}$ (ou $\rho_f$ et $\rho_f^{-1}$) ne soient pas conjugués (cf icosaèdre).

\uline{Systèmes de Schwartz} \uline{canoniques} et \uline{anticanoniques} \unsure{associé|associés} au polyèdre régulier orienté.

Quand a-t-on

\[
\Sigma_{r_1}=\Sigma_{r_2}
\]

($r_1$ et $r_2$ \deleted{directs)} repères donnés) -- disons avec $=$ orientation, i.e.

\[
\rho_{s_1}=\rho_{s_2},\inserted{\text{et}}\ \rho_{f_1}=\rho_{f_2},\quad \rho_{a_1}=\rho_{a_2}
\]

\deleted{Il semble} Posons $r_2=g\cdot r_1$ ($g\in\mathrm{Aut}\,\Pi$ bien déterminé) alors

\[
\Sigma_{r_2}=g\cdot\Sigma_{r_1}=
(\mathrm{int}(g)\rho,\mathrm{int}\,g\rho',\mathrm{int}(g)\rho'')
\]

si $\Sigma_{r_1}=(\rho,\rho',\rho'')$. Donc $\Sigma_{r_2}=\Sigma_{r_1}$ \unsure{sss|ssi} $g$ appartient au \deleted{\unsure{centre}} \inserted{commutant} $\unsure{Z|3|\mathfrak Z}$ de $\mathrm{Aut}^{+}\Pi$ \inserted{dans $\mathrm{Aut}\,\Pi$}. Dans les cas où $\mathrm{Aut}\,\Pi=Z\cdot\mathrm{Aut}^{\unsure{+|\epsilon}}(\Pi)$ \inserted{(cas des II et III)}, on voit donc que \deleted{\unsure{$\Sigma_{r_1}$ et $\Sigma_{r}$}} tous les $\Sigma_r$ conjugués par $\mathrm{Aut}^{+}(\Pi)$ (pour les $r$ directs et inverses)

% archive sheet number: 48
